Sigui $A=\begin{pmatrix}1&0\\3&1\end{pmatrix}$.

\begin{apartats}

\apartat[1,25]{0,75}
Calcula $A^2$ i $A^3$, i dedueix-ne $A^n$ per a qualsevol nombre natural $n$.

\begin{solucio}
\[
A^2=\begin{pmatrix}1&0\\3&1\end{pmatrix}\begin{pmatrix}1&0\\3&1\end{pmatrix}=\begin{pmatrix}1&0\\6&1\end{pmatrix},\qquad
A^3=A^2A=\begin{pmatrix}1&0\\9&1\end{pmatrix}.
\]
Cada vegada que es multiplica per $A$, l'element de la fila 2 i la columna 1 augmenta en 3, i la resta no
canvia: $A^n=\begin{pmatrix}1&0\\3n&1\end{pmatrix}$.
\end{solucio}

\begin{tria}{us-potencia}
\itemtria{expressio}{1}{1,25}
Calcula $A^{10}-2A^5$.

\begin{solucio}
Amb la fórmula de $A^n$:
\[
A^{10}-2A^5=\begin{pmatrix}1&0\\30&1\end{pmatrix}-2\begin{pmatrix}1&0\\15&1\end{pmatrix}
=\begin{pmatrix}1-2&0\\30-30&1-2\end{pmatrix}=\begin{pmatrix}-1&0\\0&-1\end{pmatrix}=-I.
\]
\end{solucio}

\itemtria{troba-n}{1}{1,25}
Hi ha algun nombre natural $n$ tal que $A^n=\begin{pmatrix}1&0\\60&1\end{pmatrix}$? I tal que
$A^n=\begin{pmatrix}1&0\\50&1\end{pmatrix}$? Justifica-ho.

\begin{solucio}
$A^n=\begin{pmatrix}1&0\\3n&1\end{pmatrix}$, i per tant cal que l'element de la fila 2 i la columna 1
sigui $3n$.\\
Per a la primera, $3n=60$ dona $n=20$: $A^{20}=\begin{pmatrix}1&0\\60&1\end{pmatrix}$.\\
Per a la segona, $3n=50$ dona $n=\frac{50}3$, que no és natural: cap potència de $A$ és igual a
aquesta matriu.
\end{solucio}
\end{tria}

\begin{nomesllarg}
\apartat{0,75}
Comprova que la fórmula de $A^n$ és coherent: calcula $A^n\cdot A$ fent-la servir i veu que dona la
fórmula de $A^{n+1}$.

\begin{solucio}
\[
A^n\cdot A=\begin{pmatrix}1&0\\3n&1\end{pmatrix}\begin{pmatrix}1&0\\3&1\end{pmatrix}
=\begin{pmatrix}1&0\\3n+3&1\end{pmatrix}=\begin{pmatrix}1&0\\3(n+1)&1\end{pmatrix}=A^{n+1}.
\]
\end{solucio}
\end{nomesllarg}

\end{apartats}
